\documentclass[11pt]{article}
\usepackage{amsmath}
\usepackage{amssymb}
\usepackage{geometry}
\usepackage{booktabs}
\usepackage[hidelinks]{hyperref}
\geometry{a4paper, margin=1in}

\title{O(1) Deterministic GF(4) Tensor Mapping and Sparse Residual\\
Extraction for Allele-Conditioned Computational Chemistry}
\author{Shane \and Xuening Wu}
\date{July 2026}

\begin{document}
\maketitle

\begin{abstract}
We formalize a falsifiable, iteration-free computational-chemistry program
over \(GF(4)\): nucleotide 48-mers are mapped to finite-field tensors;
mutant alleles induce a sparse residual \(\delta\); and molecular candidate
tensors \(x\) are constrained by the syndrome equation \(Hx=B\delta\), with
operators \((H,B)\) synthesized from local genomic background. The identity
\(T(s)=s+\alpha\) is deprecated as tautological. This ledger encodes algebraic
conditioning only (\texttt{clinical\_grade=false}); it does not claim affinity,
steric clearance, synthesis, or clinical performance.
\end{abstract}

\section{Introduction}
To establish a zero-trust, \(O(1)\) verification environment for allele-aware
oncology tensors, we deprecate the tautological exact-form null identity
\(T(s)=s+\alpha\). Instead, we implement a non-trivial, falsifiable constrained
syndrome equation
\begin{equation}
  Hx = B\delta
\end{equation}
to isolate spatial geometric variance derived strictly from mutant alleles,
ensuring deterministic topological constraints without biological iteration.

\section{GF(4) Isomorphism}
We establish a strict isomorphism mapping nucleotide bases to the finite field
\(GF(4)\cong\mathbb{F}_2[x]/(x^2+x+1)\) with characteristic~2
(where \(y\oplus y=0\) and \(\alpha^2+\alpha+1=0\)). The physical base levels
are mapped as follows:
\begin{equation}
\phi(A)=0,\quad \phi(C)=1,\quad \phi(G)=\alpha,\quad \phi(U)=\alpha^2.
\end{equation}

\section{Sparse Residual \(\delta\) Extraction}
For the 48-mer tensor \(s\in GF(4)^{48}\) centered at codon \(k\)
(e.g., KRAS G12), we isolate the exact spatial disparity.

First, we define the local reference tensor for the wild-type (GGU):
\begin{equation}
(s_{\mathrm{ref}}^{(k)},\, s_{\mathrm{ref}}^{(k+1)},\, s_{\mathrm{ref}}^{(k+2)})
= (\alpha,\,\alpha,\,\alpha^2).
\end{equation}

Next, we define the local mutant tensors for G12C (UGU), G12D (GAU), and
G12V (GUU) in \(O(1)\) time:
\begin{align}
\text{G12C (UGU)}
&\implies
(s_{G12C}^{(k)},\, s_{G12C}^{(k+1)},\, s_{G12C}^{(k+2)})
= (\alpha^2,\,\alpha,\,\alpha^2), \\
\text{G12D (GAU)}
&\implies
(s_{G12D}^{(k)},\, s_{G12D}^{(k+1)},\, s_{G12D}^{(k+2)})
= (\alpha,\,0,\,\alpha^2), \\
\text{G12V (GUU)}
&\implies
(s_{G12V}^{(k)},\, s_{G12V}^{(k+1)},\, s_{G12V}^{(k+2)})
= (\alpha,\,\alpha^2,\,\alpha^2).
\end{align}

We compute the pure spatial deviation tensor \(\delta\in GF(4)^{48}\) via
field addition (XOR on the bit lift):
\begin{equation}
\delta = s_{\mathrm{mut}} \oplus s_{\mathrm{ref}}.
\end{equation}

By construction, all unmutated background loci annihilate instantaneously
(\(s_{\mathrm{mut}}^{(i)}\oplus s_{\mathrm{ref}}^{(i)}=0\)). The sparse non-zero
topological residuals strictly resolve to absolute field elements:
\begin{align}
\delta_{G12C}^{(k)} &= \alpha^2 \oplus \alpha = 1, \\
\delta_{G12D}^{(k+1)} &= 0 \oplus \alpha = \alpha, \\
\delta_{G12V}^{(k+1)} &= \alpha^2 \oplus \alpha = 1.
\end{align}

\section{Non-Trivial Inverse Solver}
The sparse residuals \(\delta_{G12C}\), \(\delta_{G12D}\), and \(\delta_{G12V}\)
represent strictly orthogonal input excitation vectors. We substitute the
exact-form logic gate with the rigorous syndrome solver:
\begin{equation}
Hx = B\delta.
\end{equation}
Where \(H\) represents the topological constraint matrix, \(B\) is the boundary
mapping operator, and \(x\) is the target solution tensor representing a
candidate molecular entity under algebraic conditioning. A geometry is treated
as algebraically admissible if and only if its tensor \(x\) deterministically
absorbs the topological tension \(\delta\) generated by the allele variance.
The solution \(x\) is conditioning evidence, not a unique chemical inverse.

\section{Deprecation of Tautological Complements}
The formulation \(\mathbf{S}_{\mathrm{path}}\oplus\mathbf{S}_{\mathrm{thera}}=\mathbf{1}\)
(and its associated exact-form null \(T(s)=s+\alpha\)) is mathematically
trivial: it functions as a fixed topological schedule rather than a
falsifiable chemical reverse-solution. The mandated transition is therefore to
the non-trivial syndrome solver \(Hx=B\delta\), ensuring the solution is
physically constrained by allele variance, not axiomatically forced.

\section{Background-Conditioned Dynamic Operators}
To remediate the pathology in which identical topological tensors \(\delta\)
yielded indiscriminate solution vectors \(x\), we abandon static operator
placeholders in favor of background-conditioned dynamic operators.

\subsection{Resolution of \(\delta\)-Collision via Toeplitz Modulation}
The solution space \(x\) for \(Hx=B\delta\) is not solely a function of the
mutation tensor \(\delta\), but is fundamentally constrained by the local
genomic context \(\mathbf{S}_{\mathrm{bg}}\in GF(4)^{47}\). Operator synthesis
is defined by
\begin{equation}
    H_{\mathrm{target}} = \mathcal{F}(\mathbf{S}_{\mathrm{bg}}),\qquad
    B_{\mathrm{target}} = \mathcal{G}(\mathbf{S}_{\mathrm{bg}}),
\end{equation}
where \(\mathcal{F}\) and \(\mathcal{G}\) act as Toeplitz-based spatial filters
mapping the 47-mer background residues into the operator matrices. Thus even
when two alleles (e.g., KRAS G12C and FLT3 D835Y) generate identical sparse
residuals \(\delta[23]=1\), the induced operator \(H_{\mathrm{target}}\)
diverts the solution into non-colliding null spaces.

\subsection{Precision Probe Pipeline}
Every probe (48-mer) is anchored to its provenance ledger via SHA-256 hashes
for reproducibility of the input tensors:
\begin{itemize}
    \item \textbf{Index-23 Anchor:} the mutation site is strictly pinned at the
    median index of the 48-mer window, ensuring algebraic alignment.
    \item \textbf{Symbolic XOR Resolution:} the spatial disparity \(\delta\) is
    computed via field addition \(\delta=s_{\mathrm{mut}}\oplus s_{\mathrm{ref}}\)
    in \(GF(4)\), isolating the non-zero topological residual to the mutation
    site.
\end{itemize}

\section{Truth Boundaries}
\begin{center}
\begin{tabular}{@{}lp{0.62\textwidth}@{}}
\toprule
Claim in this ledger & Status\\
\midrule
\(GF(4)\) base isomorphism & fixed algebraic convention\\
Sparse allele residual \(\delta\) & falsifiable, allele-specific\\
Syndrome \(Hx=B\delta\) & constrained algebraic program\\
Toeplitz background operators & collision remediation for \(x\)\\
Unique molecule inverse from \(x\) & \textbf{not claimed}\\
Docking / affinity / steric clearance & \textbf{not claimed}\\
Clinical efficacy or safety & \textbf{not claimed}
(\texttt{clinical\_grade=false})\\
\bottomrule
\end{tabular}
\end{center}

Algebraic collision pathology is addressed at the tensor level; transition to
3D pose generation remains a separate geometric observer lane and is outside
the scope of this document.

\section*{Disclaimer}
This document is a computational-chemistry method ledger only. It encodes
finite-field allele isolation and background-conditioned syndrome solving. It
does not claim therapeutic effect, manufacturing readiness, or empirical
binding.

\end{document}
